\section{Experimental Setup}
\label{sec:setup}

\subsection{Models}

We benchmark seven generation models spanning dense transformers, mixture-of-experts architectures, and local, cloud-hosted, and API deployments.

\begin{itemize}
  \item \textbf{Qwen3-32B} --- 32B dense, 4-bit AWQ via Ollama on RTX 5090.
  \item \textbf{Qwen3.5-397B-A17B} --- MoE, 397B total / 17B active~\cite{qwen2025qwen3}, vLLM on H100.
  \item \textbf{GPT-OSS-20B} --- 20B open-source MoE (5B active)~\cite{gptoss2024}, Ollama on DGX Spark.
  \item \textbf{Qwen3.5-think} --- reasoning variant with chain-of-thought traces, vLLM on H100.
  \item \textbf{GPT-4o-mini} --- OpenAI API, upper-bound reference.
  \item \textbf{DeepSeek-V4-Pro} --- DeepSeek API, similarity retrieval only.
  \item \textbf{Qwen3.6-27B} --- 27B, NVFP4 quantization, vLLM on RTX 5090.
\end{itemize}

The RAGAS critic is Qwen3.5-397B-A17B in non-thinking mode on the same H100 cluster.

\subsection{Embedding and Dataset}

All configurations use \texttt{nomic-embed-text:latest}~\cite{nomic2024nomic} (768-dim, MTEB-competitive~\cite{muennighoff2022mteb}) via Ollama. We sample 50 questions per configuration from SQuAD~\cite{rajpurkar2016squad}. The shared context corpus is loaded once and reused, so retrieval differences stem solely from chunking and embedding parameters.

\subsection{Configuration Space}

Table~\ref{tab:config-space} summarizes the parameters swept in our experiments. The cartesian product of the seven sweepable parameters yields a nominal space of $M \cdot C \cdot S \cdot O \cdot R \cdot T = 224$ parameter combinations. After accounting for two structural constraints---DeepSeek-V4-Pro runs similarity retrieval only (API restriction), and semantic chunking renders chunk size/overlap inert---182 unique Phase~1 configurations produced full metric sets suitable for cross-configuration comparison. The remaining MLflow entries correspond to partial reruns, index-only stages, and failed runs excluded from analysis.

\begin{table}[t]
\centering
\caption{Configuration space for Phase~1 experiments. The cartesian product of the top seven rows yields 224 nominal combinations; 182 produce complete evaluation bundles analysed in Section~\ref{sec:results}.}
\label{tab:config-space}
\begin{tabular}{ll}
\toprule
\textbf{Parameter} & \textbf{Values} \\
\midrule
LLM Models ($M=7$) & Qwen3-32B, Qwen3.5-397B, GPT-OSS-20B, \\
           & Qwen3.5-think, GPT-4o-mini, DeepSeek-V4-Pro, Qwen3.6-27B \\
Chunking Strategies ($C=2$) & Recursive, Semantic \\
Chunk Sizes ($S=2$) & 500, 1000 \\
Chunk Overlaps ($O=2$) & 100, 200 \\
Retrieval Methods ($R=2$) & Similarity, MMR \\
Prompt Templates ($T=2$) & Concise, Detailed \\
Retrieval Top-$k$ & 12 \\
Reranker / HyDE & Disabled \\
\midrule
Nominal combinations & $M \cdot C \cdot S \cdot O \cdot R \cdot T = 224$ \\
Analysed runs (full metric set) & 182 \\
\bottomrule
\end{tabular}
\end{table}

\subsection{Prompt Templates}

Figure~\ref{fig:prompt-templates} reproduces the two system prompts. The \textbf{Concise} template enforces a single-sentence factual answer with three WikiQA-style few-shot exemplars (two shown). The \textbf{Detailed} template removes the exemplars and instructs the model to produce a full-sentence answer that includes supporting reasoning when applicable. Both share the same human-turn format (\texttt{Context:\,\ldots Question:\,\ldots Answer:}).

\begin{figure}[t]
\centering
\begin{minipage}{0.48\textwidth}
\small\textbf{Concise}\\[2pt]
\begin{lstlisting}[style=promptstyle]
Answer the question using ONLY the
provided context.
RULES:
- Give a short, direct, factual answer
  -- one sentence maximum.
- Do NOT explain reasoning, add context,
  or use markdown.
- Do NOT say 'Based on the context'.
- If the context does not contain the
  answer, say 'I cannot answer from the
  provided context.'
Examples:
Q: What circuit court is Maryland?
A: The Circuit Courts of Maryland are
   the state trial courts of general
   jurisdiction.
Q: Who wrote the music for Star Wars?
A: John Williams.
\end{lstlisting}
\end{minipage}\hfill
\begin{minipage}{0.48\textwidth}
\small\textbf{Detailed}\\[2pt]
\begin{lstlisting}[style=promptstyle]
Answer the question based only on the
provided context.
Provide a clear, complete answer in full
sentences.
Include relevant calculations or reasoning
steps when applicable.
Do not start your answer with phrases
like 'Based on the provided context' or
'Based on the given context'.
Answer directly.

Human turn (both templates):
Context: {context}

Question: {question}

Answer:
\end{lstlisting}
\end{minipage}
\caption{The two prompt templates, reproduced from \texttt{benchmark/prompt\_templates/}. Concise imposes a one-sentence factual constraint with WikiQA-style few-shot exemplars; Detailed removes exemplars and asks for full-sentence answers with reasoning.}
\label{fig:prompt-templates}
\end{figure}

\subsection{Evaluation Protocol}

RAGAS metrics (faithfulness, context recall, semantic similarity) are computed using Qwen3.5-397B-A17B as the critic model with a maximum of 4096 critic tokens per evaluation call. IR metrics (Hit@k, nDCG@k, Recall@k for $k \in \{1, 3, 5\}$) are computed against gold document labels provided by the SQuAD dataset. NLG metrics use RoBERTa-large for BERTScore~\cite{zhang2020bertscore} and standard implementations for METEOR~\cite{banerjee2005meteor}, ROUGE-L~\cite{lin2004rouge}, and BLEU~\cite{papineni2002bleu}. All metrics are reported as means over the 50-question sample, with per-question values retained for the bootstrap and significance analyses reported in Section~\ref{sec:results}.

\subsection{Reproducibility}

All code, configs, prompt templates, and MLflow bundles (model/dependency versions and git hashes) are at \url{https://github.com/niclashart/Benchmarking-Framework}. The 182 configurations can be re-executed via the \texttt{BENCHMARK\_STAGE} environment variable and the included \texttt{.env} file.
